\section{函数调用、栈和 ABI：程序如何变成可恢复的控制流}

操作系统的进程、线程、系统调用和异常处理，都依赖一个更基础的事实：程序运行不是一团代码随便跳，而是有调用约定、栈帧、返回地址、寄存器保存规则和二进制接口。先把普通函数调用讲清楚，再看系统调用和上下文切换，会少很多误解。

\subsection{为什么需要调用约定}

函数调用要解决几个问题：参数放在哪里，返回值放在哪里，被调用函数能改哪些寄存器，返回地址如何保存，栈如何对齐，谁负责清理参数。调用者和被调用者必须遵守同一套规则，否则分别编译的代码无法互相调用。这套规则叫 calling convention，是 ABI 的一部分。

\begin{longtable}{p{0.22\linewidth}p{0.34\linewidth}p{0.32\linewidth}}
\toprule
问题 & 调用约定回答 & 错误后果 \\
\midrule
参数传递 & 前几个参数在寄存器，更多参数在栈 & 函数读错输入 \\
返回值 & 使用指定寄存器或内存返回区 & 调用者拿错结果 \\
寄存器保存 & caller-saved / callee-saved & 返回后变量被破坏 \\
栈对齐 & 调用前满足固定对齐 & SIMD/函数序言/ABI 失败 \\
可变参数 & 额外寄存器和栈布局规则 & \code{printf} 一类函数解析错误 \\
\bottomrule
\end{longtable}

系统调用 ABI 和函数调用 ABI 相似但边界更硬。函数调用跨的是代码组织边界；系统调用跨的是用户/内核特权边界。

\subsection{System V x86-64 参数分类：不只是前六个整数寄存器}

常见说法是“前六个整数参数放寄存器”，这只覆盖最简单情况。System V x86-64 ABI 会按参数类别把值放进整数寄存器、SSE 寄存器或内存。整数和指针走 GPR，浮点数常走 XMM，过大或复杂聚合类型可能通过内存传递。

\begin{longtable}{p{0.22\linewidth}p{0.34\linewidth}p{0.32\linewidth}}
\toprule
参数类型 & 常见位置 & 注意点 \\
\midrule
整数/指针 & \code{rdi,rsi,rdx,rcx,r8,r9} & 更小整数会扩展或只使用低位寄存器 \\
浮点/double & \code{xmm0-xmm7} & 和整数寄存器计数分开 \\
小结构体 & 可能拆到 GPR/XMM & 取决于字段分类和大小 \\
大结构体参数 & 内存或指针传递 & 调用者要准备对象地址 \\
大结构体返回 & structure-return 隐藏指针 & 调用者传入返回缓冲区地址 \\
\bottomrule
\end{longtable}

这解释了为什么同一个 C 函数签名稍微变一下，汇编会完全不同。OS 读用户态寄存器时也要区分 ABI：普通函数 ABI 使用 \code{rcx} 作为第 4 个整数参数，而 Linux x86-64 syscall ABI 使用 \code{r10} 传第 4 个参数，因为 \code{syscall} 硬件会使用 \code{rcx} 保存返回 RIP。

\begin{lstlisting}
System V function call integer args:
  arg0 rdi, arg1 rsi, arg2 rdx, arg3 rcx, arg4 r8, arg5 r9

Linux x86-64 syscall args:
  nr rax, arg0 rdi, arg1 rsi, arg2 rdx, arg3 r10, arg4 r8, arg5 r9
\end{lstlisting}

如果把这两套 ABI 混用，系统调用会读错参数。用户态 syscall wrapper 的价值之一，就是把 C 函数调用 ABI 翻译成 syscall ABI。

\subsection{一个函数调用的机器级过程}

假设调用 \code{f(1, 2)}。伪汇编可能是：

\begin{lstlisting}
mov edi, 1      ; arg0, System V x86-64 integer arg 1
mov esi, 2      ; arg1, System V x86-64 integer arg 2
call f
; return value in eax
\end{lstlisting}

\asmcode{call f} 至少做两件事：把返回地址压入当前 \code{rsp} 指向的栈，把 \code{rip} 改成 \code{f} 的入口。返回时 \code{ret} 从栈顶取回返回地址放回 \code{rip}。

\begin{lstlisting}
call target:
  push next_rip
  rip = target

ret:
  pop rip
\end{lstlisting}

这解释了为什么栈损坏会导致程序跳飞。返回地址只是栈上的数据；若越界写覆盖返回地址，\code{ret} 就会把错误值装入 \code{rip}。现代系统用栈 canary、ASLR、NX、CFI 等机制降低这类攻击风险。

\subsection{栈帧里有什么}

一个典型栈帧可能包含返回地址、旧帧指针、保存的 callee-saved 寄存器、局部变量、临时 spill 槽和对齐填充。

\begin{lstlisting}
high address
  caller frame
  argument spill area
  return address
  old frame pointer
  saved registers
  local variables
  temporary spill slots
low address
\end{lstlisting}

优化编译器可能省略帧指针，把局部变量放寄存器，甚至内联函数让栈帧消失。但异常展开、调试器栈回溯、perf 采样和崩溃报告都需要某种方式恢复调用链。DWARF CFI 和 frame pointer 就是为了解决“如何从当前寄存器和栈恢复上一帧”。

\subsection{caller-saved 和 callee-saved}

寄存器保存规则是一种成本分摊。caller-saved 表示调用者若还需要该寄存器，就在调用前自己保存；callee-saved 表示被调用者若要使用，就必须在返回前恢复。

\begin{longtable}{p{0.22\linewidth}p{0.34\linewidth}p{0.32\linewidth}}
\toprule
类别 & 谁保存 & 适合用途 \\
\midrule
caller-saved & 调用者 & 短生命周期临时值、参数寄存器 \\
callee-saved & 被调用者 & 跨函数调用仍要保留的局部状态 \\
special registers & ABI 或硬件固定 & RSP、RIP、RFLAGS、线程指针等 \\
\bottomrule
\end{longtable}

上下文切换和函数调用不同。函数调用只遵守 ABI 保存一部分寄存器；中断或线程切换必须保存足以恢复执行的架构状态。系统调用入口还要保存用户传参寄存器，因为它们是系统调用语义的一部分。

\subsection{栈增长、栈溢出和 guard page}

很多系统的用户栈向低地址增长。函数调用越深，\code{rsp} 越小。若递归过深或局部数组太大，栈会撞到未映射区域。OS 通常在栈边界放 guard page：访问它时触发 page fault，内核判断是可增长栈还是栈溢出。

\begin{lstlisting}
stack layout:
  mapped stack pages
  guard page  // not present
  other mappings
\end{lstlisting}

内核栈更危险。用户栈溢出通常杀进程；内核栈溢出可能破坏内核对象，导致安全漏洞或 panic。因此内核栈较小、禁止大局部数组、使用专门栈处理 NMI/double fault 等关键异常。

\subsection{线程为什么需要自己的栈}

线程的“执行到哪里”不只在 \code{rip}，还在栈里。若两个线程共用同一个栈，它们的函数调用、返回地址和局部变量会互相覆盖。每个线程必须有独立用户栈；进入内核后，也需要对应的内核栈或受控异常栈保存内核路径。

\begin{lstlisting}
Thread A:
  user stack A
  kernel stack A

Thread B:
  user stack B
  kernel stack B
\end{lstlisting}

协程和用户态线程可以共享 OS 线程，但它们仍需要各自保存栈或保存状态机 continuation。无论实现方式如何，“下一步在哪里”和“局部状态在哪里”必须有地方保存。

\subsection{系统调用栈切换与函数调用的差别}

普通函数调用继续使用同一栈和同一特权级。系统调用从用户态进入内核态，不能继续信任用户栈。内核入口要切换到当前线程的内核栈，把用户寄存器保存成 \code{pt\_regs}，再调用 C 层系统调用处理函数。

\begin{lstlisting}
user:
  rax = syscall_number
  rdi = arg0
  syscall

entry:
  switch rsp from user_rsp to kernel_rsp
  push user register state into pt_regs
  call syscall_dispatch
\end{lstlisting}

这说明系统调用开销来自多层：CPU 特权切换、栈切换、寄存器保存、权限检查、对象查找、可能的调度和返回路径工作。它不是普通函数调用慢一点而已。

\subsection{red zone、栈对齐和内核代码为什么要小心}

System V x86-64 用户态 ABI 定义了 red zone：\code{rsp} 以下 128 字节在异步信号处理前不会被破坏，叶子函数可以临时使用而不调整 \code{rsp}。这对用户态优化有用，但内核或中断上下文不能盲目依赖用户态 red zone。内核入口会切换到内核栈，内核自身通常禁用或避免 red zone，因为中断可能在同一栈上保存状态。

\begin{lstlisting}
leaf function may use red zone in user ABI:
  temporary data at [rsp - 8], [rsp - 16], ...
  no sub rsp needed

kernel entry:
  do not trust user stack
  switch to kernel stack
  save pt_regs in controlled memory
\end{lstlisting}

栈对齐同样不是格式要求。许多 SIMD 指令、函数序言和 ABI 约定假设调用边界满足 16 字节对齐。若手写汇编破坏对齐，程序可能在看似无关的库函数里崩溃。内核构造初始用户栈、信号栈和 exec 栈时，必须满足用户态 ABI，否则新程序还没进入 \code{main} 就可能出错。

\subsection{varargs 为什么需要 ABI 细则}

\code{printf} 这类可变参数函数不知道调用者传了多少个参数、每个参数是什么类型，只能依赖格式串和 ABI 布局去取。System V x86-64 中，可变参数函数需要知道哪些寄存器参数被使用，必要时把寄存器保存到 register save area，\code{va\_list} 再按规则取出下一个参数。

\begin{lstlisting}
va_arg conceptual path:
  if next argument class still has register slot:
    read from saved register area
  else:
    read from overflow area on stack
  advance offset according to type and alignment
\end{lstlisting}

这说明 ABI 不是“编译器内部细节”。动态库、JIT、FFI、系统调用封装、eBPF helper、内核用户态边界都依赖调用约定。如果边界两侧对参数分类理解不同，错误不会表现为类型检查失败，而是直接读错寄存器或栈槽。

\subsection{异常展开和 C++ 异常为什么依赖 ABI}

C++ 异常抛出时，运行时要沿调用栈查找 handler，并在经过每个栈帧时执行析构。它必须知道每个函数的栈布局、返回地址、保存寄存器和清理代码位置。这些信息由 ABI、DWARF unwind table 和编译器共同维护。

\begin{lstlisting}
throw:
  capture exception object
  inspect current frame unwind info
  find landing pad or caller frame
  run destructors while unwinding
  transfer control to catch handler
\end{lstlisting}

内核通常不使用 C++ 异常，但它有类似需求：oops/panic 打印栈，perf 采样回溯调用链，用户态 signal handler 返回要恢复上下文，调试器要 unwind。所有这些都要求机器状态和栈布局可解释。

\subsection{signal frame：内核怎样临时改写用户控制流}

信号不是在用户程序旁边开一条神秘线程。内核要让当前线程在返回用户态时先执行 signal handler，于是会在用户栈或 alternate signal stack 上构造一个 signal frame，保存原来的用户上下文，并把返回路径安排到 signal trampoline。handler 结束后，\code{sigreturn} 系统调用让内核从 signal frame 恢复原上下文。

\begin{lstlisting}
deliver_signal:
  choose user stack or altstack
  write signal frame with saved registers and mask
  set user RIP = handler
  set user RSP = signal_frame
  arrange return path to sigreturn trampoline

sigreturn:
  validate frame
  restore saved user registers, RIP, RSP, RFLAGS, signal mask
\end{lstlisting}

这条路径把 ABI、栈、系统调用和安全检查连在一起。用户栈坏了，信号投递可能失败；signal frame 校验不严，用户就可能伪造上下文获得非法状态；异步信号打断普通代码，也解释了为什么可重入和 async-signal-safe 函数很重要。

\subsection{exec 如何构造第一帧用户栈}

\code{execve} 装载新程序后，用户程序还没有调用任何函数，但它需要一个初始栈，里面放 \code{argc}、\code{argv}、\code{envp}、auxv 和字符串数据。动态链接器和 libc 启动代码从这个栈读取参数，再调用 \code{main}。

\begin{lstlisting}
initial user stack:
  argc
  argv[0] pointer
  argv[1] pointer
  ...
  NULL
  envp[0] pointer
  ...
  NULL
  auxv entries
  strings for argv/envp
\end{lstlisting}

这说明进程启动也依赖 ABI。内核不能随便把参数放某个内核对象里；用户态启动代码按 ABI 从栈读取。若栈对齐或 auxv 错，程序可能在进入 \code{main} 前就崩溃。

\subsection{本节验收}

\begin{rubricbox}
\begin{itemize}
\tightlist
\item 能解释调用约定为什么必须规定参数、返回值、寄存器保存和栈对齐。
\item 能画出 \code{CALL/RET} 如何用栈保存返回地址。
\item 能说明栈帧中返回地址、旧帧指针、保存寄存器、局部变量的位置和作用。
\item 能区分函数调用保存寄存器、系统调用保存 trap frame、线程切换保存 context。
\item 能解释为什么系统调用必须切到内核栈。
\item 能解释 \code{execve} 为什么要构造初始用户栈。
\end{itemize}
\end{rubricbox}
